%% database.tex
%% Database and OGL entrypoint commands for ttrpg-core.

%% \generateDatabase[mode]{path}{UserMacro}
%% Scans path for .json files, writes .ttrpg_cache/, and \inputs the manifest.
%% Special case: when UserMacro is "LaTex" (case-insensitive), scanner reads
%% .tex files directly from path and emits \input/\include lines without JSON
%% parsing or cache generation.
%%
%% mode (default: hybrid):
%%   hybrid  — \include per chapter folder, \input per item; supports per-chapter \includeonly
%%   include — \include per item; supports per-item \includeonly
%%   input   — flat \input list; simplest, no \includeonly support
%%
%% Inside UserMacro, access values with \get{key} and \ifhas{key}{code}.
%% Fast iteration: \includeonly{.ttrpg_cache/_chap--monsters--humanoids}
\newcommand{\generateDatabase}[3][hybrid]{%
  \directlua{
    if not ttrpg_scanner then
    ttrpg_scanner = ttrpg_load_module("scanner")
    end

    local manifest = ttrpg_scanner.generate(
    "\luaescapestring{#2}",
    "\luaescapestring{#1}",
    "\luaescapestring{#3}"
    )
    tex.print("\\input{" .. manifest .. "}")
  }%
}

%% \generateDatabaseWithCallback[mode]{path}{UserMacro}{LuaCallbackName}
%% Same as \generateDatabase, but allows a Lua chapter callback to emit custom
%% heading/chapter preamble per top-level folder.
\newcommand{\generateDatabaseWithCallback}[4][hybrid]{%
\directlua{
if not ttrpg_scanner then
ttrpg_scanner = ttrpg_load_module("scanner")
end

local chapter_cb = "\luaescapestring{#4}"
local wrapper_cb = chapter_cb .. "_wrapper"
local wrapper_opt = nil
if type(_G[wrapper_cb]) == "function" then
wrapper_opt = wrapper_cb
end

local manifest = ttrpg_scanner.generate(
"\luaescapestring{#2}",
"\luaescapestring{#1}",
"\luaescapestring{#3}",
{ chapter_callback = chapter_cb, chapter_wrapper_callback = wrapper_opt }
)
tex.print("\\input{" .. manifest .. "}")
}%
}

%% \importDir[mode]{path}
%% Defaults to include mode so each imported file starts on a new page.
\NewDocumentCommand{\importDir}{O{include} m}{%
  \generateDatabase[#1]{#2}{LaTeX}%
}

%% \importDirWithCallback[mode]{path}{LuaCallbackName}
\NewDocumentCommand{\importDirWithCallback}{O{include} m m}{%
  \generateDatabaseWithCallback[#1]{#2}{LaTeX}{#3}%
}

%% Shared recursive folder counters.
%% The generic form accepts a comma-separated folder list and a format.
%% Use \\countTex and \\countJson for common cases.
\NewDocumentCommand{\countFiles}{O{tex} m}{%
  \directlua{
    if not ttrpg_folder_counter then
    ttrpg_folder_counter = ttrpg_load_module("folder_counter")
    end
    tex.print(tostring(ttrpg_folder_counter.count(
    "\luaescapestring{#2}",
    "\luaescapestring{#1}"
    )))
  }%
}

\NewDocumentCommand{\countTex}{m}{%
  \countFiles[tex]{#1}%
}

\NewDocumentCommand{\countJson}{m}{%
  \countFiles[json]{#1}%
}

%% \generateOGL{csvpath}
%% Reads a CSV of copyright credits and typesets a full OGL chapter.
%%
%% CSV columns: title, year, publisher, authors
%% An optional header row is skipped if the first field is "title".
%% Fields containing commas must be quoted with double-quotes.
%% LaTeX macros in any field (e.g. \textbf{}, \&) are passed through as-is.
\newcommand{\generateOGL}[1][./credits.csv]{%
  \directlua{
    if not ttrpg_ogl then
    local ok, mod = pcall(ttrpg_load_module, "ogl")
    if not ok then
    error("ttrpg-core: failed to load ogl.lua: " .. tostring(mod))
    end
    ttrpg_ogl = mod
    end

    ttrpg_ogl.generate("\luaescapestring{#1}")
  }%
}

%% \RandomTable{csvpath}
%% Reads a CSV and typesets a random-roll table, with a "dN" die column
%% auto-generated (N = number of data rows, numbered 1..N) — including the
%% header, e.g. 10 data rows produces a "d10" column heading.
%%
%% CSV columns: first row = header names for the data columns (the die
%% column itself is not a CSV column). Each remaining row is one roll
%% result. The first data column stretches to fill the table width; any
%% further columns are centered — matching the common "d10 | Name | Qty |
%% Value" table shape.
\newcommand{\RandomTable}[1]{%
  \directlua{
    if not ttrpg_randomtable then
    local ok, mod = pcall(ttrpg_load_module, "randomtable")
    if not ok then
    error("ttrpg-core: failed to load randomtable.lua: " .. tostring(mod))
    end
    ttrpg_randomtable = mod
    end

    ttrpg_randomtable.generate("\luaescapestring{#1}")
  }%
}
